% ============================================================================
% Victor Barres — Curriculum Vitae
% ============================================================================
%
% Build (from this directory):
%   latexmk -xelatex barres_cv.tex
%
% Or one-shot:
%   xelatex barres_cv.tex
%
% Then copy/publish the resulting PDF to ../assets/pdf/barres_cv.pdf
% (or run ../bin/build_cv.sh).
% ============================================================================

\documentclass[10.5pt,a4paper]{article}

% ---- Page geometry ---------------------------------------------------------
\usepackage[a4paper,left=0.7in,right=0.7in,top=0.6in,bottom=0.6in]{geometry}
\usepackage[parfill]{parskip}

% ---- Fonts -----------------------------------------------------------------
\usepackage{fontspec}
\setmainfont{Palatino}                % macOS system Palatino (.ttc)
\usepackage{microtype}

% ---- Color + links ---------------------------------------------------------
\usepackage[dvipsnames,svgnames]{xcolor}
\definecolor{accent}{HTML}{1f3a5f}   % deep slate-blue, understated
\definecolor{muted}{HTML}{555555}

\usepackage[hidelinks,colorlinks=true,urlcolor=accent,linkcolor=accent]{hyperref}

% ---- Typography helpers ----------------------------------------------------
\usepackage{needspace}
\usepackage{titlesec}
\titleformat{\section}
  {\needspace{6\baselineskip}\large\bfseries\color{accent}}
  {}{0pt}{}
  [\vspace{-3pt}\color{accent}\hrule height 0.5pt\vspace{2pt}]
\titlespacing*{\section}{0pt}{10pt}{4pt}

% Compact lists
\usepackage{enumitem}
\setlist[itemize]{leftmargin=1.1em, topsep=2pt, parsep=0pt, itemsep=2pt, label=\textcolor{accent}{\textbullet}}

% Header band
\newcommand{\name}[1]{{\Huge\bfseries\color{accent}#1}}
\newcommand{\role}[1]{{\large\color{muted}#1}}

% Role / org entry
% \position{employer}{location}{start--end}
%          {title}{notes (optional)}
\newcommand{\position}[5]{%
  \noindent
  \textbf{#1}\hfill\textcolor{muted}{#3}\\
  \textit{#4}\ifx&#5&\else, #5\fi\hfill\textcolor{muted}{#2}\par
  \vspace{2pt}
}

% Education entry: same shape, slightly compressed
\newcommand{\edu}[5]{%
  \noindent
  \textbf{#1}\hfill\textcolor{muted}{#3}\\
  \textit{#4}\ifx&#5&\else, #5\fi\hfill\textcolor{muted}{#2}\par
  \vspace{2pt}
}

% Inline reference to a paper in Selected Publications
\newcommand{\paper}[3]{%
  \noindent
  \textbf{#1}\par
  \textcolor{muted}{\small #2.\hfill #3}\par
  \vspace{4pt}
}

\pagestyle{empty}

% ============================================================================
\begin{document}
% ============================================================================

% ---- Header ----------------------------------------------------------------
\begin{minipage}{\textwidth}
\name{Victor Barres, PhD}\\[2pt]
\role{Research Scientist \textperiodcentered\ Conversational Agents \& Evaluation}\\[4pt]
\textcolor{muted}{%
  \href{mailto:victor.barres@gmail.com}{victor.barres@gmail.com}
  \ \textperiodcentered\
  \href{https://victorbarres.github.io}{victorbarres.github.io}
  \ \textperiodcentered\
  \href{https://www.linkedin.com/in/victor-barres/}{linkedin.com/in/victor-barres}
  \ \textperiodcentered\
  \href{https://scholar.google.com/citations?user=RpPIx5AAAAAJ}{Google Scholar}%
}
\end{minipage}

% ============================================================================
\section*{Experience}
% ============================================================================

\position{Mercor}{San Francisco, CA}{Jun 2026 – Present}
  {Research Scientist}{}
\begin{itemize}
  \item Founding member of the \textbf{APEX research team} (\href{https://www.mercor.com/blog/why-mercor-is-building-a-research-team/}{AI Productivity Index}), measuring how capable AI is at economically valuable work.
  \item Research on conversational AI agents inside organizations --- how they work with people, how they are transforming work itself, and how to measure that well enough to understand it.
  \item Building a bridge between organization theory and AI evaluation with Copenhagen Business School: \emph{Harder Than the Task} (2026), a taxonomy of organizational difficulty applied to 522 AI benchmarks.
\end{itemize}

\position{Sierra}{San Francisco, CA}{Sep 2024 – May 2026}
  {AI Researcher}{}
\begin{itemize}
  \item Led the \textbf{$\tau$-Bench family} of agent benchmarks --- research direction, codebase, public leaderboard (\href{https://taubench.com}{taubench.com}), and ecosystem of extensions: $\tau^2$-Bench (dual control), $\tau$-Knowledge (knowledge retrieval), $\tau$-Voice (full-duplex voice), and $\tau^3$-Bench (the combined release with community-contributed task and code improvements). Three papers from the family accepted to \textbf{ICML 2026}, with $\tau^2$-Bench as an \textbf{oral}. Continued contributing to the line after leaving: Hyper-$\tau$-Bench (agents that build agents), $\tau$-i18n, and $\tau$-Entity (2026).
  \item Kickstarted Sierra's internal \textbf{agent evaluation system}: set up the proof of concept and did the original work on user-simulation-driven stress testing for pre-release agent hardening across deployed customer-service domains.
\end{itemize}

\position{Elemental Cognition}{Remote}{Jun 2023 – Jul 2024}
  {Senior Researcher (NLP)}{}
\begin{itemize}
  \item Developed core algorithms for EC's hybrid AI architecture, integrating LLMs with formal reasoning engines.
  \item Led the integration of multi-round RAG with formal models of causation for knowledge-intensive multi-hop QA, published as \emph{From Generating Answers to Building Explanations} at NAACL 2025 (Industry).
  \item Built conversational agents for two reasoning-heavy settings: a logic-programming assistant that elicits and translates user requirements into a logic program, and an ASP-reasoner-backed constraint-solving agent.
\end{itemize}

\position{Uniphore.ai}{Palo Alto, CA}{Dec 2019 – Jun 2023}
  {Principal NLP Scientist \& NLP Research Lead}{previously Senior NLP Scientist (2019–2021)}
\begin{itemize}
  \item Built and led Uniphore's NLP Lab (7+ scientists). Set short-, mid-, and long-term research strategy for understanding human–human spoken conversation in the contact-center domain.
  \item Shipped foundational and applied work across LLM benchmarking, prompt engineering, structured knowledge grounding, summarization, semantic parsing, dialogue-act and dialogue-flow modeling, and real-time caller-intent recognition.
  \item Kickstarted the team, deployed first NER / redaction / call-driver models, and stood up cloud infrastructure and CI/CD for production NLP.
\end{itemize}

\position{Astound.ai}{Menlo Park, CA}{Sep 2017 – Dec 2019}
  {AI Scientist (NLP)}{}
\begin{itemize}
  \item Core member of the team that built, from the ground up, Astound's domain-expert conversational agent for IT / HR / Finance enterprise self-service.
  \item Led the neuro-symbolic side of the agent's NLP stack — word/sentence embeddings, semantic parsing, knowledge-graph embeddings, and semantic inference — integrated with a modular dialogue manager (domain classification, intent recognition, search, small talk, sub-flows).
  \item Built NLP-side APIs over Astound's ontology and extended it to support inference-driven multi-turn clarification, establishing a separation of concerns between knowledge, lexicon, and grammar engineers.
\end{itemize}

% ============================================================================
\section*{Service}
% ============================================================================

\begin{itemize}
  \item Co-organizer, \href{https://usersim-workshop.github.io/}{Grounded User Simulation for Model Evaluation and Training}, NeurIPS 2026 Workshop (Paris, Dec 2026).
  \item Mercor representative, \href{https://github.com/openenv/OpenEnv/blob/main/GOVERNANCE.md}{OpenEnv Technical Committee} (2026--present).
  \item Talks: \href{https://victorbarres.github.io/assets/html/fit-for-purpose-training-agents-irl-2026.html}{Fit for Purpose} (Training Agents IRL, Oct 2026, \href{https://youtu.be/CKFUx7COIwQ?t=1170}{video}); \href{https://victorbarres.github.io/assets/html/tau2-bench-icml2026-oral.html}{$\tau^2$-Bench} (ICML 2026 Oral, Jul 2026).
  \item Organizer \& judge, Sierra $\tau^2$-Bench Custom Track, \href{https://rdi.berkeley.edu/agentx-agentbeats.html}{AgentX–AgentBeats} competition (Berkeley RDI, 2025–2026).
\end{itemize}

% ============================================================================
\section*{Education}
% ============================================================================

\edu{University of Southern California}{Los Angeles, CA}{Sep 2010 – Aug 2017}
  {PhD, Computational Cognitive Neuroscience}{}
Dissertation: \emph{Schema Architecture for Language–Vision Interactions: A Computational Cognitive Neuroscience Model of Language Use.} Advisor: Michael A. Arbib.
\vspace{4pt}

\edu{École Normale Supérieure}{Paris, France}{Jun 2008 – Jun 2010}
  {MS, Cognitive Science}{}
\vspace{2pt}

\edu{École Polytechnique}{Palaiseau, France}{Jun 2003 – Jun 2006}
  {MS, Physics}{}
\vspace{2pt}

\edu{Lycée Louis-le-Grand}{Paris, France}{Jun 2001 – Jun 2003}
  {Classes Préparatoires – Mathematics \& Physics}{}

% ============================================================================
\section*{Languages}
% ============================================================================

\textbf{English} (fluent) \textperiodcentered\ \textbf{French} (native) \textperiodcentered\ \textbf{Spanish} (conversational)

% ============================================================================
\section*{Selected Publications}
% ============================================================================

\paper{Harder Than the Task: From Technical to Organizational Difficulty in AI Benchmarking}
  {V. Barres, A. Barton, N. J. Foss, H. M. Nielsen}
  {arXiv:2611.03494, 2026}

\paper{$\tau^\tau$-Bench (Hyper-$\tau$-Bench): An Environment for End-To-End, Realistic Agent Construction}
  {Q. Shi, K. Dhandhania, K. Narasimhan, V. Barres}
  {arXiv:2609.04611, 2026}

\paper{$\tau$-i18n ($\tau$-Multilingual): Benchmarking Voice Agents Across Languages}
  {S. Ray, E. dos Santos Paiva, R. Valenzuela, K. Narasimhan, K. Dhandhania, V. Barres}
  {arXiv:2609.35820, 2026}

\paper{$\tau$-Voice: Benchmarking Full-Duplex Voice Agents on Real-World Domains}
  {V. Barres*, S. Ray*, K. Dhandhania*, K. Narasimhan}
  {arXiv:2603.13686, 2026. Accepted to ICML 2026.}

\paper{$\tau$-Knowledge: Evaluating Conversational Agents over Unstructured Knowledge}
  {Q. Shi, A. Zytek, P. Razavi, K. Narasimhan, V. Barres}
  {arXiv:2603.04370, 2026. Accepted to ICML 2026.}

\paper{$\tau^2$-Bench: Evaluating Conversational Agents in a Dual-Control Environment}
  {V. Barres*, H. Dong*, S. Ray, X. Si, K. Narasimhan}
  {arXiv:2506.07982, 2025. ICML 2026 (Oral).}

\paper{From Generating Answers to Building Explanations: Integrating Multi-Round RAG and Causal Modeling for Scientific QA}
  {V. Barres, C. McFate, A. Kalyanpur, K. K. Saravanakumar, L. Moon, N. Seifu, A. Bautista-Castillo}
  {NAACL 2025 (Industry)}

\paper{LLM-ARC: Enhancing LLMs with an Automated Reasoning Critic}
  {A. Kalyanpur, K. Saravanakumar, V. Barres, J. Chu-Carroll, D. Melville, D. Ferrucci}
  {arXiv:2406.17663, 2024}

\paper{Template Construction Grammar: From Visual Scene Description to Language Comprehension and Agrammatism}
  {V. Barres, J. Lee}
  {Neuroinformatics 12(1), 2014}

\vspace{4pt}
\textcolor{muted}{\small * Equal contribution. Full list: \href{https://victorbarres.github.io/publications/}{victorbarres.github.io/publications/}}

\end{document}
